\subsection{Gorenstein 环与平坦代数}

命题 12. 设 \(\rho: A \to B\) 为 Noether 局部环之间的一个局部同态，它使 B 成为平坦 A-模。下列条件等价：

\begin{enumerate}
\def\labelenumi{(\roman{enumi})}
\item
  B 是 Gorenstein 环；
\item
  A 与 \(\kappa_{\Lambda} \otimes_{\Lambda} B\) 都是 Gorenstein 环。
\end{enumerate}

我们先处理环 A 与 B 都是 Artin 环的情形。用 C 表示局部环 \(\kappa_{\Lambda} \otimes_{\mathrm{A}} \mathrm{B}\)；其剩余域 \(\kappa_{\mathrm{C}}\) 等同于 \(\kappa_{\mathrm{B}}\)。由于 B 在 A 上平坦，B-模 \(\operatorname{Hom}_{\mathrm{B}}(\mathrm{C}, \mathrm{B})\) 同构于 \(\operatorname{Hom}_{\mathrm{A}}(\kappa_{\mathrm{A}}, \mathrm{A}) \otimes_{\mathrm{A}} \mathrm{B}\) (I, § 2, 第 10 节, 命题 11)，从而同构于 \(\operatorname{Hom}_{\mathrm{A}}(\kappa_{\mathrm{A}}, \mathrm{A}) \otimes_{\kappa_{\Lambda}} \mathrm{C}\)。由此得出一列同构：

\[
\begin{array}{c} \operatorname{Hom} _ {\mathrm{B}} (\kappa_ {\mathrm{B}}, \mathrm{B}) \longrightarrow \operatorname{Hom} _ {\mathrm{C}} (\kappa_ {\mathrm{C}}, \operatorname{Hom} _ {\mathrm{B}} (\mathrm{C}, \mathrm{B})) \dashrightarrow \operatorname{Hom} _ {\mathrm{C}} (\kappa_ {\mathrm{C}}, \operatorname{Hom} _ {\Lambda} (\kappa_ {\mathrm{A}}, \mathrm{A}) \otimes_ {\kappa_ {\mathrm{A}}} \mathrm{C}) \\ \longrightarrow \operatorname{Hom} _ {\Lambda} (\kappa_ {\Lambda}, \mathrm{A}) \otimes_ {\kappa_ {\mathrm{A}}} \operatorname{Hom} _ {\mathrm{C}} (\kappa_ {\mathrm{C}}, \mathrm{C}). \end{array}
\]

特别地，有 \(\left[\mathrm{Hom}_{\mathrm{B}}(\kappa_{\mathrm{B}},\mathrm{B}):\kappa_{\mathrm{B}}\right] = \left[\mathrm{Hom}_{\mathrm{A}}(\kappa_{\mathrm{A}},\mathrm{A}):\kappa_{\mathrm{A}}\right]\left[\mathrm{Hom}_{\mathrm{C}}(\kappa_{\mathrm{C}},\mathrm{C}):\kappa_{\mathrm{C}}\right]\)，于是命题由第 7 节的引理 1 得出。

转到一般情形。如果在命题陈述中把“Gorenstein”一词换为“Macaulay”，则本命题是 § 2, 第 7 节命题 10 的一个特殊情形。于是可设环 A、B 与 C = \(\kappa_{\mathrm{A}} \otimes_{\mathrm{A}} \mathrm{B}\) 都是 Macaulay 的。B-模 C 是 Macaulay 的 (§ 2, 第 1 节, 例 4)。令 \(r = \dim(\mathrm{A})\)，\(s = \dim(\mathrm{C})\)。存在一个 A-正则序列 ( \(x_1, \ldots, x_r\) )，其元素属于 \(\mathfrak{m}_{\mathrm{A}}\)，以及一个序列 ( \(y_1, \ldots, y_s\) )，其元素属于 \(\mathfrak{m}_{\mathrm{B}}\) 且对 B-模 C 正则 (§ 2, 第 3 节, 命题 3)；用 \(x\) 表示前者生成的 A 的理想，用 \(\mathfrak{y}\) 表示后者生成的 B 的理想。序列 ( \(y_1, \ldots, y_s, \rho(x_1), \ldots, \rho(x_r)\) ) 是 B-正则的 (§ 1, 第 6 节, 命题 11)，且 A-模 B/ \(\mathfrak{y}\) 是平坦的（同上引文, 命题 10）。由 A/ \(x\) 到 B/(xB + \(\mathfrak{y}\) ) 的、由 \(\rho\) 出发取商得到的同态因此使 B/(xB + \(\mathfrak{y}\) ) 成为平坦的 (A/x)-模，且环 \(\kappa_{\mathrm{A}/x} \otimes_{\mathrm{A}/x} \mathrm{B}/(xB + \mathfrak{y})\) 同构于 C/ \(\mathfrak{y}\) C。于是，考虑到第 7 节的例 3，命题由证明的第一部分得出。

推论 1. -- 设 \(\mathfrak{p}: A \to B\) 为使 B 成为平坦 A-模的 Noether 环同态。下列条件等价：

\begin{enumerate}
\def\labelenumi{(\roman{enumi})}
\item
  B 是 Gorenstein 环；
\item
  （相应地 (iii)）对 B 的每个素理想（相应地极大理想）\(\mathfrak{q}\)，环 \(\mathrm{A}_{\mathfrak{p}^{-1}(\mathfrak{q})}\) 与 \(\kappa (\mathfrak{p}^{-1}(\mathfrak{q}))\otimes_{\mathrm{A}}\mathrm{B}\) 都是 Gorenstein 的。
\end{enumerate}

此外，若 B 是忠实平坦的 \(\Lambda\) -模，则这些条件蕴含 A 是 Gorenstein 环。

设 q 为 B 的一个素理想；令 \(\mathfrak{p} = \rho^{-1}(\mathfrak{q})\) 。环 \(B_{q}\) 同构于 \(B_{p}\) 的一个分式环，在 \(A_{p}\) 上平坦。为使 \(B_{q}\) 是 Gorenstein 环，必须且只需环 \(A_{p}\) 与 \(\kappa(\mathfrak{p}) \otimes_{A_{\mathfrak{p}}} B_{\mathfrak{q}}\) 都是 Gorenstein 环 (命题 12)。

\((\mathrm{i}) \Rightarrow (\mathrm{ii})\)：设 \(\mathfrak{q}\) 为 B 的一个素理想，并令 \(\mathfrak{p} = \rho^{-1}(\mathfrak{q})\)。若 B 是 Gorenstein 环，则 \(B_{\mathfrak{q}}\) 亦然；于是由前所证，\(A_{\mathfrak{p}}\) 与 \(\kappa(\mathfrak{p}) \otimes_{A_{\mathfrak{p}}} B_{\mathfrak{q}}\) 都是 Gorenstein 的。由 § 2, 第 6 节的注，\(\kappa(\mathfrak{p}) \otimes_{\mathrm{A}} B\) 在每个素理想处的局部化这时都是 Gorenstein 环，这蕴含 \(\kappa(\mathfrak{p}) \otimes_{\mathrm{A}} B\) 是 Gorenstein 环，从而得到 (ii)。

\begin{enumerate}
\def\labelenumi{(\roman{enumi})}
\setcounter{enumi}{1}
\tightlist
\item
  \(\Rightarrow\) (iii)：这是显然的。
\end{enumerate}

(iii)⇒(i)：对 B 的每个极大理想 n，由证明的开头（取 q = n 应用）可知 \(B_{n}\) 是 Gorenstein 环，从而得到 (i)。

若 B 是忠实平坦 A-模，则映射 \(^{q}\mathfrak{p}:\operatorname{Spec}(B)\longrightarrow\operatorname{Spec}(\Lambda)\) 是满射 (II, § 2, 第 5 节, 命题 11 的推论 4)，由此得到最后一个断言。

推论 2. 设 A 为 Noether 环，J 为 A 的理想，Â 为 A 关于 J-进拓扑的分离完备化。考察下列条件：

\begin{enumerate}
\def\labelenumi{(\roman{enumi})}
\item
  A 是 Gorenstein 环；
\item
  \(\widehat{\mathrm{A}}\) 是 Gorenstein 环；
\item
  对 A 的每个包含 J 的极大理想 m，\(A_{m}\) 是 Gorenstein 环；
\item
  对 A 的每个满足 \(p + J \neq A\) 的素理想 p，\(A_{p}\) 与 \(\kappa(\mathfrak{p}) \otimes_{A} \widehat{A}\) 都是 Gorenstein 环。
\end{enumerate}

条件 (ii) 至 (iv) 等价，且均被 (i) 蕴含。当理想 J 含于 A 的根中时，条件 (i) 至 (iv) 等价。

我们由推论 1 得出本推论，其方法与由 § 2, 第 7 节的命题 9 得出其推论 4 的方法相同：只需在证明中把“Macaulay”一词换为“Gorenstein”。

推论 3.--- 设 A 为 Gorenstein 环，\((\mathrm{T}_{i})_{i\in\mathrm{I}}\) 为有限的未定元族。环 \(\mathrm{A}[(\mathrm{T}_{i})_{i\in\mathrm{I}}]\) 与 \(\mathrm{A}[((\mathrm{T}_{i})_{i\in\mathrm{I}}]]\) 都是 Gorenstein 环。

对每个域 \(k\)，环 \(k[\mathrm{T}]\) 都是 Gorenstein 的 (第 7 节, 例 5)；此外，环 A{[}T{]} 是 Noether 的。因此它是 Gorenstein 环 (推论 1)。对 \(\operatorname{Card}(\mathrm{I})\) 用归纳法即知 \(\mathrm{A}[(\mathrm{T}_i)_{i\in \mathbf{I}}]\) 是 Gorenstein 环；然后应用推论 2。

